% Pulsar threat model -- included by pulsar.tex.
% Mirrors docs/threat-model.md (now deleted; this file is the
% canonical LaTeX form).

\subsection{Adversary classes}

We give an adversary the union of the following capabilities,
parameterized by $t$ (the corruption threshold) and the corruption
model:

\paragraph{Static corruption (default).}

\begin{itemize}[leftmargin=*, itemsep=0pt]
\item Adversary $\mathcal{A}$ selects up to $t-1$ parties to corrupt
\textbf{before} protocol start. From start to end of the protocol,
those $t-1$ parties are controlled; the remaining honest parties run
the protocol faithfully.
\item $\mathcal{A}$ sees all messages on the broadcast channel + all
messages on the authenticated point-to-point channels into corrupted
parties.
\item $\mathcal{A}$ learns each corrupted party's full state (secret
share, randomness, per-round values) including session-state values.
\item $\mathcal{A}$ can selectively delay messages within the round
window but cannot forge messages from honest parties on
authenticated channels.
\end{itemize}

\paragraph{Adaptive corruption (stronger model).}

\begin{itemize}[leftmargin=*, itemsep=0pt]
\item Adversary $\mathcal{A}$ can corrupt up to $t-1$ parties at any
time during the protocol, including mid-round, conditioned on the
transcript so far.
\item Once corrupted, a party remains corrupted for the rest of the
protocol.
\item All other capabilities as in static.
\end{itemize}

The Pulsar security analysis (\texttt{spec/security-games.tex})
provides the precise game definitions. Static security is the
baseline; adaptive security is a stretch goal for the first MPTC
submission and a hard requirement for the second-round response.

\paragraph{Mobile corruption (proactive resharing).}
For long-lived deployments where the validator set rotates between
epochs (Pulsar's \texttt{reshare/} direction), the relevant model is
HJKY97 mobile adversary: $\mathcal{A}$ can corrupt up to $t-1$
parties \textbf{per epoch}, with the corruption set varying across
epochs. The adversary may \emph{uncorrupt} a party (forget what it
learned) at epoch boundaries.

Pulsar proactive resharing inherits the soundness argument from
Pulsar's \texttt{reshare/} (with the constant-time fix from F9 of
the Hanzo HIP-0077 red review). The MPTC submission documents the
mobile-adversary game explicitly.

\subsection{Out-of-scope adversary capabilities}

\begin{itemize}[leftmargin=*, itemsep=0pt]
\item Compromise of the cryptographic hash function (SHAKE-256 /
cSHAKE-256 / KMAC-256). NIST FIPS 202 + SP 800-185 are taken as
building blocks.
\item Side-channel attacks on the \emph{physical} hardware running a
party (power analysis, EM, fault injection). Mitigation discussed in
\S\ref{sec:knownlimits} as future work.
\item Compromise of the BIP-32/BIP-39 mnemonic from which
device-specific keys are derived. The HD-derivation security
argument lives in HIP-0077 \S``Identity''; the Pulsar layer
assumes parties' input keys are uniformly distributed.
\end{itemize}

\subsection{In-scope security goals}

\paragraph{Threshold unforgeability (TS-UF).}
For any PPT adversary $\mathcal{A}$ controlling up to $t-1$ parties,
$\mathcal{A}$ cannot produce a valid Pulsar signature $\sigma^*$
on any message $\mu^*$ that an honest signing ceremony has not
produced.

\paragraph{Robustness / identifiable abort.}
If a corrupted party deviates from the protocol, the honest parties:

\begin{enumerate}[leftmargin=*, itemsep=0pt]
\item Detect the deviation.
\item Identify the deviating party.
\item Produce signed evidence of the deviation suitable for slashing
/ reputation reduction.
\end{enumerate}

\paragraph{Output indistinguishability from FIPS 204.}
A Pulsar signature $\sigma$ on message $\mu$ against group public
key $\mathit{pk}$ is computationally indistinguishable from a
FIPS~204 ML-DSA signature on the same $(\mathit{pk}, \mu)$.
Specifically: an adversary that can distinguish the two with
non-negligible advantage breaks the underlying M-LWE assumption.

\paragraph{Bias resistance.}
The randomness used in each party's contribution (Gaussian samples,
challenge polynomials) is statistically indistinguishable from
uniform even given the adversary's view of corrupted parties'
randomness.

\paragraph{DKG soundness.}
Pedersen DKG over $\Rq^k$ produces a public key $\mathit{pk}$
distributed uniformly over the valid M-LWE public-key space, with no
party holding more than $1/n$ partial information about the
corresponding secret key.

\subsection{Assumed channels}

\begin{center}
\begin{tabular}{lll}
\toprule
channel & property & purpose \\
\midrule
Broadcast & authenticated, eventual-delivery & round transcripts,
  public coins \\
Point-to-point & authenticated, encrypted & private share / blind
  delivery \\
Public bulletin & authenticated, append-only & DKG commitments,
  reshare records \\
\bottomrule
\end{tabular}
\end{center}

The reference implementation provides these channels via TCP + TLS.
Production deployments (e.g.\ Lux consensus + Quasar) bind these to
their existing P2P transport.

\subsection{Trusted parameters}

\begin{itemize}[leftmargin=*, itemsep=0pt]
\item The cyclotomic ring $\Rq = \Zq[X] / (X^{256} + 1)$ and
parameters $(q, k, \ell, \eta, \beta, \omega, \ldots)$ exactly as in
FIPS~204 ML-DSA-65.
\item The hash family: SHAKE-256 (FIPS 202), cSHAKE-256 / KMAC-256 /
TupleHash-256 (NIST SP 800-185). Customization tags are pinned in
\texttt{spec/pulsar.tex}.
\item The Gaussian sampler parameters from FIPS~204 \S3.3 / \S4.2.
\end{itemize}

No trusted setup beyond NIST primitives.

\subsection{Failure modes addressed}

The MPTC submission documents how Pulsar handles each failure
mode that broke earlier threshold-lattice schemes:

\begin{itemize}[leftmargin=*, itemsep=0pt]
\item \textbf{Rejection-sampling restart leaking secret information
across rounds}: Pulsar's PRNG-key per-round domain separation
(CRIT-1 fix from upstream's 2026-05-03 amendment) ports forward.
\item \textbf{Constant-time share comparison}:
\texttt{dkg2.constTimePolyEqual} pattern ported from upstream
Pulsar (\texttt{reshare/commit.go} F9 fix).
\item \textbf{Information leaks via logging}: zero stdlib
\texttt{log} use in any secret-touching path; all logging via
\texttt{luxfi/log} at fields that are explicitly public.
\item \textbf{Deterministic quorum selection in resharing leaking
adversary positioning}: Pulsar uses a beacon-derived per-reshare
quorum selection (HIP-0077 \S``reshare quorum'' F10 fix).
\item \textbf{Cross-mode hash-suite confusion}: NIST profile is
exclusively SHA-3 family; no BLAKE3 in Pulsar ever.
\end{itemize}
